% !TEX root = main.tex
\section{Collatz Character and Collatz Kernel}~\label{sec:collatz-character-and-collatz-kernel}

Although the periodicity of reduce vectors constitutes a significant result, we cannot stop there. As we learned in the previous section, not all reduce vectors lead to positive integers; indeed, many reduce vectors eventually resolve into fractions as the depth of the backtracking increases. Consequently, we are more interested in determining the conditions under which a number can be traced back to an integer, as well as the distribution of the resulting values.

To deal with this, we have to introduce two operators.
\begin{definition}~\label{def:collatz-character-and-collatz-kernel}
Given $b \in\widehat{\mathbb{N}}_\beta$, define
\[ \chr b = \min\{v\in\mathbb{N}: \alpha\mid(\beta^{v}b-\gamma) \text{ and } \beta\nmid(\beta^{v}b-\gamma)/\alpha \} \]
called the Collatz character of $b$. If the set is empty, we say $b$ is a primitive source and define $\chr b = \infty$.
If $b$ is not primitive source, we define
\[ \cor b = \beta^{\chr b}b \]
as the source core of $b$.
Given $B \subset \widehat{\mathbb{N}}_\beta$, we denote
\begin{align*}
    \chr B &= \{ \chr b \mid b\in B - \mathbb{S}\}, \\
    \cor B &= \{ \cor b \mid b\in B - \mathbb{S}\}.
\end{align*}
\end{definition}
\begin{definition}
We define
\[\mathbb{S}(B) = B\cup\bigcup_{b\in B}\{ a\in \widehat{\mathbb{N}}_{\beta} \mid a\rightsquigarrow b \}\]
as all numbers that will fall into to $B$.
When $B = \{b\}$, that's $B$ contains only one number, we also denote it as $\mathbb{S}(b)$.
\end{definition}
\begin{proposition}
    $\cor b\equiv \gamma \pmod{\alpha}$.
\end{proposition}
And according properties of the discrete logarithm, we have
\begin{proposition}~\label{prop:chr-le-ord}
    If $\chr a < \infty$, $\chr a \le \ord_{\alpha}(\beta)$.
\end{proposition}
\begin{proposition}~\label{prop:collatz-character-finite-generated}
    Denote $\Gamma = [\gamma G(\beta;\alpha)]/\alpha$. We ahve
    \[ \chr a <\infty \Longleftrightarrow a\fmod \alpha \in \Gamma. \]
    Here $G(\beta; \alpha)$ follows \cref{def:generated-subgroup}.
\end{proposition}
\begin{proof}
    If $\chr a <\infty$, we can consider $\cor a$. Let $v = \ord_{\alpha}\beta - \chr a$,  we have
    \[ \cor a = \beta^{\chr a}a \equiv\beta^{\ord_{\alpha}\beta - v}a \equiv \beta^{-v}a \equiv\gamma. \pmod{\alpha} \]
    Thus $a\equiv \gamma\cdot\beta^{v} \pmod{\alpha}$ and $v\le \ord_{\alpha}\beta$. Thus $a\in[\gamma G(\beta; \alpha)]/\alpha$.

    If $a\in[\gamma G(\beta; \alpha)]/\alpha$, then $\exists v\le \ord_{\alpha}\beta$, such that $a\equiv \gamma\cdot\beta^{v} \pmod{\alpha}$. Then let $u = \ord_{\alpha}\beta - v \le \ord_{\alpha}\beta$, we have $\beta^{u}a\equiv\gamma \pmod{\alpha}$, which equivalents to $\alpha \mid (\beta^{u}a-\gamma)$. If $u > 0$, then $\beta\nmid(\beta^{u}a-\gamma)$, since $\gcd(\beta, \alpha) = 1$, thus we have $\beta\nmid (\beta^{u}a-\gamma)/\alpha$. If $\beta\mid(\beta^{u}a-\gamma)$, since $\gcd(\beta, \gamma) = 1$, we have $u = 0$. Thus $a-\gamma \equiv 0 \pmod{\alpha}$. Then we have $\beta\nmid (\beta^{\ord_{\alpha}\beta}a-\gamma) \equiv a-\gamma \equiv 0 \pmod{\alpha}$. Anyway, $\chr a\le \ord_{\alpha}\beta < \infty$.
\end{proof}
\begin{definition}
    We denote $\overline{\Gamma} = \Gamma \cup\{\infty\}$.
    And we define $\mathbb{S}(b)_{i} = \{ a\rightsquigarrow b \mid \chr a = i \}$, $\mathbb{S}_{i} = \{ b\in\widehat{\mathbb{N}}_{\beta} \mid \chr b = i \}$.
\end{definition}
With these, we have
\begin{proposition}
    $\mathbb{S}(b) = \bigcup_{i\in\overline{\Gamma}}\mathbb{S}(b)_{i}$. $\widehat{\mathbb{N}}_{\beta} = \bigcup_{i\in\overline{\Gamma}}\mathbb{S}_{i}$.
\end{proposition}
Actually, $\mathbb{S}_{\infty}$ is exactly the set of all primitive sources. And we have
\begin{proposition}~\label{prop:chr-are-periodic}
    If $i = \chr a$, then $\mathbb{S}_{i} = \{ b\in\widehat{\mathbb{N}}_{\beta} \mid b \equiv a \pmod{\alpha} \}$.
\end{proposition}
And hence we also denote $\mathbb{S}[a] = \mathbb{S}_i$. We have
\begin{proposition}
    $\chr \mathbb{S}[a] = \{\chr a\}$.
\end{proposition}
Furthermore, from \cref{cor:modulo-alpha-of-f_1}, we have
\begin{proposition}
    If $\alpha \mid \gamma$, then $\Phi_1 = 1$, and then $\chr b = \left\{\begin{aligned}
        &1, &&\alpha\mid b, \\
        &\infty, &&\alpha\nmid b.
    \end{aligned}\right.$\newline
    And thus $\mathbb{S}_{\infty} = \{b \in \widehat{\mathbb{N}}_\beta: a\nmid b\}$.
    But if $\alpha \nmid \gamma$, then $\forall b\in\alpha\widehat{\mathbb{N}}_\beta$, $b \in \mathbb{S}_{\infty}$.
\end{proposition}

This proposition states that whether $\alpha \mid \gamma$ or $\alpha \nmid \gamma$ but $b \in \alpha\widehat{\mathbb{N}}_\beta$, these are all trivial cases. Thus when $\alpha \nmid \gamma$, given $b \in \mathbb{S}_{\infty}$. Or it is not reversible modulo $\alpha$, or
\[\Log_{\beta, \alpha}(b^{-1} \cdot \gamma) = \varnothing. \]
And $b$ is not reversible modulo $\alpha$ equivalents $\gcd(b, \alpha) > 1$. Hence we have
\begin{definition}
    Suppose $\alpha \nmid \gamma$. If $\gcd(b, \alpha) > 1$, we say they are normal primitive source.
    Others we call them abnormal primitive source.
    All normal primitive source are denoted as 
    \[\hat{\mathbb{S}}_{\infty} = \{b \in \widehat{\mathbb{N}}_{\beta} \mid \gcd(b, \alpha) > 1\} = \bigcup_{p\in\Pi(\mathbb{P}\mid\alpha)} p\widehat{\mathbb{N}}_{{\beta}}, \]
    all abnormal primitive source are denoted as
    \[ \tilde{\mathbb{S}}_{\infty} = \{b \in \widehat{\mathbb{N}}_{\beta} \mid \gcd(b, \alpha) = 1 \text{ and } \Log(b^{-1}\cdot\gamma) = \varnothing\}. \]
\end{definition}
Immidiately, we have
\begin{proposition}
    $\mathbb{S}_{\infty} = \tilde{\mathbb{S}}_{\infty} \cup \hat{\mathbb{S}}_{\infty}$.
\end{proposition}
\begin{proposition}~\label{prop:primitive-source-of-simple-collatz-system}
    If $\alpha\in\mathbb{P}$ and $\alpha\nmid\gamma$, then $\tilde{\mathbb{S}}_{\infty} = \varnothing$ and $\hat{\mathbb{S}}_{\infty} = \alpha\mathbb{N}_{\ge 1}$.
\end{proposition}
